\begin{abstract}
Natural-language requirements can be individually clear yet impose incompatible obligations when read together. This report studies a prototype that translates requirements with a language model, checks the resulting formulas with Z3, and returns review questions with witnesses and explicit assumptions. The method tests competing translations, asks about relations between separately named outcomes, and checks proposed OpenSpec changes against existing requirements.

The experiments examine conflict detection, ambiguity, change analysis and scale. In E18, outcome-pair questions increased planted conflict detections from 27/36 to 34/36 across repeated runs, while clean runs with a conflict alarm increased from 0/18 to 3/18. Second-reader experiments found more planted ambiguities and raised more questions about clean specifications. E19 tested removal handling and found 35/36 planted gaps. In E20, reusing a saved base model reduced mean formalization requests per change from 68 to 9, after an initial mean 61 requests per base; conflict detections fell from 54/54 to 50/54. Full formalization remains the default. The scale exercises reached 100 requirements, with input and representation limits still apparent.

These results support further study of assisted requirements review. They do not establish faithful translation, production defect rates or human review savings. The evidence comes mainly from project-authored examples, and the snapshot representation can confuse successive states or separate objects. Reuse currently follows scenario edits and can miss a change confined to a requirement sentence. The next work is independent evaluation, broader edit coverage and better vocabulary selection.
\end{abstract}
